% Contenido del póster: Ma, Xu, Anwar y Lu (2023), Global Finance Journal. DOI 10.1016/j.gfj.2023.100884
% Elaborado con apoyo de IA a partir del resumen oficial del artículo (OpenAlex), 2026-11-12.
\def\DOI{10.1016/j.gfj.2023.100884}
\def\TituloEs{Competencia bancaria y uso del crédito en la sombra: evidencia de las plataformas de préstamos}
\def\TituloOriginal{Banking competition and the use of shadow credit: Evidence from lending marketplaces}
\def\Autores{Ma, Xu, Anwar y Lu}
\def\Anio{2023}
\def\Revista{Global Finance Journal}

\def\Problema{Cuando la banca está concentrada, algunos prestatarios solventes pueden quedar fuera del crédito bancario y recurrir a plataformas en línea, un tipo de crédito «en la sombra».}
\def\Objetivo{Determinar si la competencia bancaria afecta la solvencia de los prestatarios que recurren al crédito en la sombra.}
\def\Metodo{Análisis de datos detallados de una de las principales plataformas de préstamos en línea de China, comparando a los prestatarios según el grado de competencia del mercado bancario de su zona.}
\def\Resultados{\begin{itemize}
  \item Los prestatarios de mercados bancarios concentrados son, en general, más solventes que los de mercados competitivos.
  \item El fenómeno es más fuerte cuando la competencia proviene de bancos nuevos y cuando los prestatarios son jóvenes o nuevos en la plataforma.
\end{itemize}}
\def\Conclusion{Los prestatarios migran de los bancos a las plataformas cuando la banca está concentrada, y una mayor competencia bancaria puede mejorar la inclusión financiera al flexibilizar la evaluación crediticia.}
\def\Relevancia{El estudio relaciona el crédito de plataforma con la estructura del mercado bancario. En el Perú, las tasas a la pequeña y microempresa son altas (la de las cajas rurales no bajó de \cifra{teaCracPymeMin}~\% entre 2022 y 2026), pero los datos públicos del financiamiento participativo no informan la solvencia de los receptores, por lo que no puede saberse si las plataformas atraen a prestatarios de mejor o peor calidad.}
\def\PalabrasClave{competencia bancaria, crédito en la sombra, plataformas de préstamos, solvencia, China}
% --- Ficha del estudio (herramientas/posteres/ficha_articulos.py) ---
\def\Ambito{China (plataforma de préstamos en línea)}
\def\Enfoque{Empírico con datos detallados de prestatarios}
\def\FiguraPeru{3}
% --- fin de la ficha ---
